\documentclass[a4paper, 12pt]{article}
\usepackage{graphicx} % Required for inserting images
\usepackage{fullpage}
\usepackage{amsmath}
\usepackage[dvipsnames]{xcolor}
\usepackage{float}
\usepackage{geometry}
\usepackage{biblatex}
\geometry{margin=1in}
\usepackage{enumitem}
\usepackage{parskip}
\usepackage{pgfplots}
\pgfplotsset{compat=1.18}
\usepackage[normalem]{ulem}
\usepackage[version=4]{mhchem}
\usepackage{tikz}
\usepackage{microtype}
\usepackage{gensymb}
\usepackage{caption}
\usepackage{cancel}
\usepackage{nicefrac}

\newcommand{\lab}[1]{\: \text{#1}}
\newcommand{\degC}{$\degree$C \,}
\newcommand{\degF}{$\degree$F \,}
\newcommand{\R}{\left(0.0821 \: \frac{L \cdot atm}{mol \cdot \text{K}}\right)}
\newcommand{\cunits}{$\frac{J}{g \degree \text{C}}$}
\newcommand{\Hf}{$\Delta H_\text{f}$} % heat of formation
\newcommand{\mathHf}{\Delta H_\text{f}} %heat of formation in math mode

\usepackage{hyperref}
\hypersetup{
colorlinks = true,
linkcolor = teal,
citecolor=teal,         % Set color for citations
filecolor=teal,         % Set color for file links
urlcolor=teal           % Set color for URLs
}


\begin{document}

\begin{center}
\textbf{{\Large AP Chem Test 1}} \\
30 September 2026
\end{center}

\section*{Atomic structure and periodicity}

\subsection*{Average atomic mass}

Average atomic mass is the summation of each mass number multiplied by its relative abundance (\%).

\textbf{Example:} Neon has three isotopes: Ne--20 (19.992 amu, 90.48\%), Ne--21 (20.994 amu, 0.27\%), and Ne--22 (21.991 amu, 9.25\%).

To find the average atomic mass:

$$19.992 \times 0.9048 + 20.994 \times 0.0027 + 21.991 \times 0.0925 = \boxed{20.180 \lab{amu}}$$

\subsection*{Electron configuration and d-orbital stabilisation}

\textbf{d-orbital stabilisation} refers to the d-orbitals preferring to be half-filled or completely filled (5 or 10) rather than slightly underfilled or overfilled (4, 6, 9).

The noble gas configuration of chromium, for example, appears to be [Ar]4s$^2$3d$^4$. However, to stabilise the d-orbital, one electron is borrowed from the $s$-orbital, making the final noble gas configuration [Ar]4s$^13$d$^5$. Similarly:

\begin{table}[H]
\begin{tabular}{l l}
Ag & [Kr]5s$^24$d$^9$ $\to$ [Kr]5s$^14$d$^{10}$ \\
Cu & [Ar]4s$^23$d$^9$ $\to$ [Ar]4s$^13$d$^{10}$ \\
Fe$^{2+}$ & [Ar]4s$^23$d$^6$ $\to$ [Ar]4s$^13$d$^5$ \\
Fe$^{3+}$ & [Ar]4s$^2$3d$^6$ $\to$ [Ar]3d$^5$
\end{tabular}
\end{table}
For ions, electrons are removed from the outermost level (4 for the two iron ions shown).

\subsection*{Photo electron spectroscopy}

\begin{figure}[H]
\centering
\includegraphics[width=0.5\textwidth]{~/Dropbox/apchemistry/test1/PES}
\end{figure}

\subsection*{Periodic table trends and Coulomb's Law}

Coulomb's Law states that

\begin{equation}
\label{eq:coulomb}
F \propto \frac{q_1q_2}{r^2}
\end{equation}

where $F$ is the force (repulsive if positive, attractive if negative), $q_1$ and $q_2$ are the respective charges, and $r$ is the distance between the two particles.

This becomes important when considering periodic table trends such as:

\begin{description}
\item[ionisation energy (IE)] energy required to remove an electron from a gaseous atom or ion
\item[electron affinity (EA)] energy released with the addition of an electron to a gaseous ion or atom
\item[electronegativity (EN)] an atom's tendency to attract electrons toward itself
\end{description}

\textbf{EA and size increase down and to the left.}

In a given period, the greater the number of protons, the smaller the atom will be. However, as you go down the periodic table, the number of electron shells increases along with size. Electron affinity increases down and to the left, but its absolute value (since more energy released = more negative) increases up and to the right.

\textbf{IE and EN increase up and to the right.}

Ionisation energy increases as you go up, since the number of electron shells decreases, making the radius ($r^2$) smaller. Across the period, the number of protons increases, pulling the electron shell closer to it, making it harder to remove an electron. The higher positive charge pulling the electrons close to the nucleus is also the reason for the trend in electronegativity.

\subsection*{Multiple ionisation energies}

The first, second, third, and fourth ionisation energies of Al are as follows:

$I_1 = 580 \lab{kJ/mol}$\\
$I_2 = 1815 \lab{kJ/mol}$\\
$I_3 = 2740 \lab{kJ/mol}$\\
$I_4 = 11600 \lab{kJ/mol}$

The jump in ionisation energy from $I_3$ to $I_4$ can be attributed to Al's position on the periodic table; as it loses three electrons, it is stripped of its outer electron shell, giving it a full octet and making it harder to remove another electron from the new valence shell.

\section*{Bonding}

\subsection*{Bond length diagram}



\section{Intermolecular forces}

\end{document}